\documentclass[12pt,a4paper]{article}
\usepackage[utf8]{inputenc}
\usepackage[T1]{fontenc}
\usepackage[english]{babel}
\usepackage{lmodern}  % Latin Modern fonts with proper accent support
\usepackage{amsmath}
\usepackage{amssymb}
\usepackage{amsthm}
\usepackage{graphicx}
\usepackage{hyperref}
\usepackage{listings}
\usepackage{xcolor}
\usepackage{geometry}
\usepackage{fancyhdr}
\usepackage{tcolorbox}
\usepackage{enumitem}
\usepackage{multicol}

\geometry{margin=2.5cm}
\pagestyle{fancy}
\fancyhf{}
\rhead{Dryad Language - Theoretical Foundation}
\lhead{\thesection}
\cfoot{\thepage}

\theoremstyle{definition}
\newtheorem{definition}{Definição}[section]
\newtheorem{theorem}{Teorema}[section]
\newtheorem{axiom}{Axioma}[section]
\newtheorem{property}{Propriedade}[section]

\title{\textbf{Dryad Programming Language} \\ 
\Large Fundamentos Teóricos e Especificação Formal}
\author{Theoretical Foundation Document \\ Version 1.0}
\date{\today}

\begin{document}

\maketitle
\newpage

\tableofcontents
\newpage

\begin{abstract}
Este documento apresenta os fundamentos teóricos completos da linguagem de programação Dryad, uma linguagem multi-paradigma de tipagem dinâmica inspirada em JavaScript/TypeScript, mas com características próprias que a distinguem. A especificação cobre aspectos lexicais, sintáticos, semânticos, sistema de tipos, modelo de execução, sistema de módulos, concorrência, e o modelo de compilação híbrido (interpretador + bytecode + AOT). Este documento é puramente teórico e não contém implementações de código, servindo como referência fundamental para compreensão conceitual da linguagem e base para futuras reimplementações.
\end{abstract}

\section{Introdução e Filosofia da Linguagem}

\subsection{Visão Geral}

Dryad é uma linguagem de programação multi-paradigma projetada com os seguintes objetivos fundamentais:

\begin{itemize}
    \item \textbf{Familiaridade Sintática}: Sintaxe inspirada em JavaScript/TypeScript para reduzir curva de aprendizado
    \item \textbf{Flexibilidade de Paradigma}: Suporte simultâneo para programação procedural, orientada a objetos e funcional
    \item \textbf{Tipagem Dinâmica Opcional}: Sistema de tipos dinâmico com anotações opcionais para documentação e verificação
    \item \textbf{Modelo de Execução Híbrido}: Suporte para interpretação direta de AST, compilação para bytecode e compilação AOT para código nativo
    \item \textbf{Concorrência Nativa}: Primitivas de concorrência integradas (threads, async/await, mutexes)
    \item \textbf{Sistema de Módulos Robusto}: Import/export moderno com suporte a módulos nativos via FFI
    \item \textbf{Gerenciamento Automático de Memória}: Garbage collection transparente
\end{itemize}

\subsection{Características Distintivas}

\begin{definition}[Características Fundamentais]
Dryad se diferencia por:
\begin{enumerate}
    \item \textbf{Arquitetura de Execução Tripla}: Interpretador de AST, VM de bytecode stack-based, e compilador AOT
    \item \textbf{Sistema de Erros Estruturado}: Todos os erros são categorizados com códigos únicos, localizações e sugestões
    \item \textbf{Operador de Namespace}: Suporte ao operador \texttt{::} para acesso a namespaces (estilo C++/Rust)
    \item \textbf{Native Directives}: Carregamento explícito de módulos nativos via diretivas \texttt{\#<module>}
    \item \textbf{Template Strings}: Interpolação de strings nativa com \texttt{`\$\{expr\}`}
\end{enumerate}
\end{definition}

\subsection{Princípios de Design}

\begin{axiom}[Princípios Fundamentais]
A linguagem Dryad é regida pelos seguintes princípios:
\begin{itemize}
    \item \textbf{Clareza sobre Performance}: Prioriza-se legibilidade e clareza do código sobre otimização prematura
    \item \textbf{Consistência Sintática}: Estruturas sintáticas seguem padrões consistentes
    \item \textbf{Erros Informativos}: Todo erro deve fornecer localização, contexto e sugestões de correção
    \item \textbf{Composabilidade}: Primitivas da linguagem devem ser combináveis
    \item \textbf{Zero Surpresas}: Comportamento da linguagem deve ser previsível e documentado
\end{itemize}
\end{axiom}

\section{Estrutura Lexical}

\subsection{Alfabeto e Caracteres}

\begin{definition}[Alfabeto Dryad]
O alfabeto da linguagem Dryad é composto por:
\begin{itemize}
    \item Caracteres Unicode UTF-8
    \item Letras: \texttt{[a-zA-Z]}
    \item Dígitos: \texttt{[0-9]}
    \item Símbolos especiais: \texttt{\{ \} ( ) [ ] ; , . : = + - * / \% \& | \^ \textasciitilde ! < > ? @ \# \$ \textbackslash ' " `}
    \item Whitespace: espaços, tabulações, quebras de linha
\end{itemize}
\end{definition}

\subsection{Tokens e Classificação}

\begin{definition}[Sistema de Tokens]
A análise lexical de Dryad produz os seguintes tipos de tokens:

\begin{enumerate}
    \item \textbf{Identificadores}: \texttt{Identifier(String)} \\
    Padrão: \texttt{[a-zA-Z\_][a-zA-Z0-9\_]*}
    
    \item \textbf{Literais Numéricos}: \texttt{Number(f64)} \\
    Suporta: decimal, hexadecimal (\texttt{0x}), binário (\texttt{0b}), octal (\texttt{0o})
    
    \item \textbf{Literais String}: \texttt{String(String)} \\
    Delimitadores: \texttt{"} ou \texttt{'}
    
    \item \textbf{Literais Booleanos}: \texttt{Boolean(bool)} \\
    Valores: \texttt{true}, \texttt{false}
    
    \item \textbf{Literal Null}: \texttt{Literal("null")}
    
    \item \textbf{Template Strings}: \\
    \texttt{TemplateStart}, \texttt{TemplateContent(String)}, \\
    \texttt{InterpolationStart}, \texttt{InterpolationEnd}, \texttt{TemplateEnd}
    
    \item \textbf{Palavras-chave}: \texttt{Keyword(String)}
    
    \item \textbf{Operadores}: \texttt{Operator(String)}
    
    \item \textbf{Símbolos}: \texttt{Symbol(char)}
    
    \item \textbf{Diretivas Nativas}: \texttt{NativeDirective(String)} \\
    Formato: \texttt{\#<module\_name>}
    
    \item \textbf{Fim de Arquivo}: \texttt{Eof}
\end{enumerate}
\end{definition}

\subsection{Palavras Reservadas}

\begin{definition}[Conjunto de Keywords]
As palavras-chave reservadas de Dryad são organizadas por categoria:

\textbf{Declaração de Variáveis:}
\begin{itemize}[noitemsep]
    \item \texttt{let} (variável mutável)
    \item \texttt{const} (variável imutável)
\end{itemize}

\textbf{Controle de Fluxo:}
\begin{itemize}[noitemsep]
    \item \texttt{if}, \texttt{else}
    \item \texttt{for}, \texttt{while}, \texttt{do}
    \item \texttt{break}, \texttt{continue}
    \item \texttt{match}, \texttt{in}
\end{itemize}

\textbf{Funções:}
\begin{itemize}[noitemsep]
    \item \texttt{function}, \texttt{fn}
    \item \texttt{async}, \texttt{await}
    \item \texttt{return}
    \item \texttt{thread}, \texttt{mutex}
\end{itemize}

\textbf{Orientação a Objetos:}
\begin{itemize}[noitemsep]
    \item \texttt{class}, \texttt{new}
    \item \texttt{extends}, \texttt{this}, \texttt{super}
    \item \texttt{static}, \texttt{public}, \texttt{private}, \texttt{protected}
\end{itemize}

\textbf{Módulos:}
\begin{itemize}[noitemsep]
    \item \texttt{import}, \texttt{export}
    \item \texttt{use}, \texttt{as}, \texttt{from}
\end{itemize}

\textbf{Tratamento de Erros:}
\begin{itemize}[noitemsep]
    \item \texttt{try}, \texttt{catch}, \texttt{finally}
    \item \texttt{throw}
\end{itemize}

\textbf{Literais:}
\begin{itemize}[noitemsep]
    \item \texttt{true}, \texttt{false}, \texttt{null}
\end{itemize}
\end{definition}

\subsection{Operadores}

\begin{definition}[Sistema de Operadores]
Dryad possui os seguintes operadores organizados por categoria:

\textbf{Operadores Aritméticos:}
\begin{itemize}[noitemsep]
    \item \texttt{+} (adição/concatenação)
    \item \texttt{-} (subtração)
    \item \texttt{*} (multiplicação)
    \item \texttt{/} (divisão)
    \item \texttt{\%} (módulo)
    \item \texttt{**} (exponenciação)
\end{itemize}

\textbf{Operadores Bitwise:}
\begin{itemize}[noitemsep]
    \item \texttt{\&} (AND bitwise)
    \item \texttt{|} (OR bitwise)
    \item \texttt{\textasciicircum}, \texttt{\textasciicircum\textasciicircum} (XOR bitwise)
    \item \texttt{<<} (shift left)
    \item \texttt{>>} (shift right aritmético)
    \item \texttt{>>>} (shift right unsigned)
    \item \texttt{<<<} (rotate left)
    \item \texttt{\%\%} (módulo bitwise)
\end{itemize}

\textbf{Operadores de Comparação:}
\begin{itemize}[noitemsep]
    \item \texttt{==} (igualdade)
    \item \texttt{!=} (desigualdade)
    \item \texttt{<} (menor que)
    \item \texttt{>} (maior que)
    \item \texttt{<=} (menor ou igual)
    \item \texttt{>=} (maior ou igual)
\end{itemize}

\textbf{Operadores Lógicos:}
\begin{itemize}[noitemsep]
    \item \texttt{\&\&} (AND lógico)
    \item \texttt{||} (OR lógico)
    \item \texttt{!} (NOT lógico)
\end{itemize}

\textbf{Operadores de Atribuição:}
\begin{itemize}[noitemsep]
    \item \texttt{=} (atribuição)
    \item \texttt{+=}, \texttt{-=}, \texttt{*=}, \texttt{/=} (atribuições compostas)
\end{itemize}

\textbf{Operadores de Incremento/Decremento:}
\begin{itemize}[noitemsep]
    \item \texttt{++} (incremento, prefix ou postfix)
    \item \texttt{--} (decremento, prefix ou postfix)
\end{itemize}

\textbf{Operadores Especiais:}
\begin{itemize}[noitemsep]
    \item \texttt{=>} (arrow function)
    \item \texttt{\#} (native directive prefix)
    \item \texttt{::} (namespace access)
\end{itemize}
\end{definition}

\subsection{Comentários e Whitespace}

\begin{definition}[Comentários]
Dryad suporta dois tipos de comentários:
\begin{itemize}
    \item \textbf{Linha}: \texttt{// comentário até o fim da linha}
    \item \textbf{Bloco}: \texttt{/* comentário multi-linha */}
\end{itemize}
Comentários são descartados durante análise lexical e não fazem parte da AST.
\end{definition}

\begin{property}[Whitespace]
Espaços em branco (espaços, tabs, newlines) são geralmente insignificantes, exceto:
\begin{itemize}
    \item Como separadores entre tokens
    \item Em strings literais
    \item Em template strings
\end{itemize}
Não há significado sintático para indentação (diferente de Python).
\end{property}

\subsection{Diretivas Nativas e Desambiguação Lexical}

\begin{definition}[Diretiva Nativa]
Uma diretiva nativa tem o formato:
\[
\texttt{\#<module\_name>}
\]
onde \texttt{module\_name} é um identificador válido delimitado por \texttt{<} e \texttt{>}.
\end{definition}

\begin{property}[Desambiguação do Símbolo \#]
O caractere \texttt{\#} em Dryad possui dois contextos de uso:
\begin{enumerate}
    \item \textbf{Diretiva Nativa}: Quando seguido imediatamente (sem whitespace) por \texttt{<}, é interpretado como início de diretiva nativa
    \item \textbf{Operador Futuro}: Reservado para potenciais operadores futuros
\end{enumerate}

O analisador lexical garante distinção clara através das seguintes regras:
\begin{itemize}
    \item Se \texttt{\#} é seguido de \texttt{<} (sem espaços): token \texttt{NativeDirective}
    \item Caso contrário: erro léxico (caractere inesperado)
\end{itemize}
\end{property}

\begin{axiom}[Princípio da Não-Ambiguidade]
O analisador lexical de Dryad garante que:
\[
\forall \text{ sequência de entrada } s, \exists! \text{ interpretação léxica } t \text{ tal que } lex(s) = t
\]
Ou seja, para qualquer entrada válida existe uma única interpretação de tokens possível.
\end{axiom}

\textbf{Exemplos de Análise}:
\begin{verbatim}
#<io>          → NativeDirective("io")
# <io>         → Erro Léxico (espaço após #)
//#<io>        → Comentário (descartado)
"#<io>"        → String Literal
\end{verbatim}

\section{Sistema de Tipos}

\subsection{Fundamentos do Sistema de Tipos}

\begin{definition}[Tipagem Dinâmica]
Dryad implementa um \textbf{sistema de tipos dinâmico} onde:
\begin{itemize}
    \item Tipos são associados a valores, não a variáveis
    \item Variáveis podem mudar de tipo durante execução
    \item Verificação de tipo ocorre em tempo de execução
    \item Anotações de tipo são opcionais e servem para documentação
\end{itemize}
\end{definition}

\subsection{Tipos Primitivos}

\begin{definition}[Tipos Primitivos de Dryad]
Os tipos primitivos fundamentais são:

\begin{enumerate}
    \item \textbf{number}: Número de ponto flutuante (IEEE 754 f64) \\
    Literais: \texttt{42}, \texttt{3.14}, \texttt{0xFF}, \texttt{1e10}
    
    \item \textbf{string}: Sequência de caracteres UTF-8 \\
    Literais: \texttt{"hello"}, \texttt{'world'}, \texttt{`template`}
    
    \item \textbf{bool}: Valor booleano \\
    Literais: \texttt{true}, \texttt{false}
    
    \item \textbf{null}: Valor nulo/indefinido \\
    Literal: \texttt{null}
    
    \item \textbf{any}: Tipo universal (implícito) \\
    Pode conter qualquer valor
\end{enumerate}
\end{definition}

\subsection{Tipos Compostos}

\begin{definition}[Arrays]
Arrays são coleções heterogêneas ordenadas de tamanho dinâmico.

\textbf{Sintaxe de Tipo}: \texttt{Type[]}

\textbf{Propriedades}:
\begin{itemize}
    \item Indexação baseada em zero
    \item Tamanho dinâmico
    \item Podem conter tipos mistos quando declarados como \texttt{any[]}
    \item Acesso por índice: \texttt{arr[i]}
\end{itemize}
\end{definition}

\begin{definition}[Tuples]
Tuplas são coleções heterogêneas de tamanho fixo.

\textbf{Sintaxe de Tipo}: \texttt{(Type1, Type2, ...)}

\textbf{Propriedades}:
\begin{itemize}
    \item Tamanho fixo definido na criação
    \item Cada posição tem tipo potencialmente diferente
    \item Acesso por índice numérico: \texttt{tuple.0}, \texttt{tuple.1}
\end{itemize}
\end{definition}

\begin{definition}[Functions]
Funções são valores de primeira classe.

\textbf{Sintaxe de Tipo}: \texttt{fn(ParamType1, ParamType2, ...) -> ReturnType}

\textbf{Propriedades}:
\begin{itemize}
    \item Podem ser atribuídas a variáveis
    \item Podem ser passadas como argumentos
    \item Podem ser retornadas de outras funções
    \item Suportam closures (capturam escopo)
\end{itemize}
\end{definition}

\begin{definition}[Objects]
Objetos são coleções de pares chave-valor.

\textbf{Propriedades}:
\begin{itemize}
    \item Chaves são strings
    \item Valores podem ser de qualquer tipo
    \item Acesso via notação ponto: \texttt{obj.prop}
    \item Acesso via notação colchetes: \texttt{obj["prop"]}
\end{itemize}
\end{definition}

\begin{definition}[Classes]
Classes são tipos definidos pelo usuário com estado e comportamento.

\textbf{Propriedades}:
\begin{itemize}
    \item Suportam herança via \texttt{extends}
    \item Possuem propriedades de instância e estáticas
    \item Possuem métodos de instância e estáticos
    \item Suportam modificadores de acesso: \texttt{public}, \texttt{private}, \texttt{protected}
\end{itemize}
\end{definition}

\subsection{Anotações de Tipo}

\begin{definition}[Sintaxe de Anotações]
Anotações de tipo aparecem após identificadores com sintaxe \texttt{:}

\textbf{Exemplos}:
\begin{itemize}[noitemsep]
    \item Variáveis: \texttt{let x: number = 42}
    \item Parâmetros: \texttt{function add(a: number, b: number)}
    \item Retorno: \texttt{function add(...): number}
\end{itemize}

\textbf{Regras}:
\begin{itemize}
    \item Se omitida, tipo é inferido ou \texttt{any}
    \item Anotações são verificadas em parse time mas não enforced em runtime
    \item Servem primariamente para documentação
\end{itemize}
\end{definition}

\subsection{Modo Estrito de Tipos (Futuro)}

\begin{definition}[Strict Type Mode]
O \textbf{modo estrito} é uma extensão futura planejada onde anotações de tipo são validadas estaticamente pelo compilador AOT para gerar código otimizado.
\end{definition}

\begin{property}[Benefícios do Modo Estrito]
Quando habilitado, o modo estrito fornece:
\begin{itemize}
    \item \textbf{Otimização AOT}: Compilador gera código de máquina especializado sem overhead de verificação dinâmica de tipos
    \item \textbf{Eliminação de Verificações Runtime}: Tipos conhecidos em tempo de compilação permitem eliminar type checks
    \item \textbf{Especialização de Código}: Funções genéricas podem ser monomorfizadas para tipos concretos
    \item \textbf{Detecção de Erros Antecipada}: Incompatibilidades de tipo detectadas antes da execução
    \item \textbf{Melhor Performance}: Redução de 30-50\% no overhead de tipo em hot paths
\end{itemize}
\end{property}

\begin{axiom}[Compatibilidade de Modo]
O modo estrito deve ser \textbf{opt-in} e compatível com código dinâmico:
\begin{itemize}
    \item Módulos podem especificar \texttt{"use strict types"} no topo do arquivo
    \item Código estrito pode interoperar com código dinâmico através de fronteiras explícitas
    \item Funções individuais podem ser marcadas como \texttt{@strict} ou \texttt{@dynamic}
    \item Conversão implícita na fronteira: tipos dinâmicos são verificados ao entrar em código estrito
\end{itemize}
\end{axiom}

\textbf{Exemplo Conceitual (Futuro)}:
\begin{verbatim}
"use strict types";

// Função estrita: tipos verificados estaticamente
function add(a: number, b: number): number {
    return a + b;  // Compilador garante que a + b é number
}

// Compilado para código de máquina sem type checks:
// mov rax, [a]
// add rax, [b]
// ret

// Função dinâmica (fallback)
@dynamic
function process(x: any): any {
    return x + 1;  // Runtime type check necessário
}
\end{verbatim}

\begin{property}[Otimizações Habilitadas]
Com tipos estaticamente conhecidos, o compilador AOT pode aplicar:
\begin{enumerate}
    \item \textbf{Inline Caching Eliminado}: Sem necessidade de cache polimórfico
    \item \textbf{Devirtualização}: Chamadas de método resolvidas estaticamente
    \item \textbf{Escape Analysis}: Alocações de stack substituem heap quando possível
    \item \textbf{Dead Code Elimination}: Branches impossíveis removidas
    \item \textbf{SIMD Vectorization}: Loops sobre arrays tipados vetorizados automaticamente
\end{enumerate}
\end{property}

\subsection{Coerção e Compatibilidade de Tipos}

\begin{property}[Conversões Implícitas]
Dryad realiza as seguintes conversões automáticas:
\begin{itemize}
    \item \texttt{number + string} $\rightarrow$ \texttt{string} (concatenação)
    \item \texttt{bool} em contexto aritmético $\rightarrow$ \texttt{1} (true) ou \texttt{0} (false)
    \item \texttt{null} em comparações tem tratamento especial
    \item Arrays e objetos são comparados por referência
\end{itemize}
\end{property}

\section{Expressões}

\subsection{Hierarquia de Precedência}

\begin{definition}[Tabela de Precedência]
Expressões em Dryad seguem a seguinte hierarquia de precedência (do maior para o menor):

\begin{enumerate}
    \item \textbf{Primárias}: \texttt{.}, \texttt{[]}, \texttt{()}, \texttt{new} (left-associative)
    \item \textbf{Postfix}: \texttt{++}, \texttt{--} (left-associative)
    \item \textbf{Prefix}: \texttt{++}, \texttt{--}, \texttt{!}, \texttt{-}, \texttt{+}, \texttt{\textasciitilde} (right-associative)
    \item \textbf{Exponenciação}: \texttt{**} (right-associative)
    \item \textbf{Multiplicativas}: \texttt{*}, \texttt{/}, \texttt{\%} (left-associative)
    \item \textbf{Aditivas}: \texttt{+}, \texttt{-} (left-associative)
    \item \textbf{Bitwise Shift}: \texttt{<<}, \texttt{>>}, \texttt{>>>}, \texttt{<<<} (left-associative)
    \item \textbf{Relacionais}: \texttt{<}, \texttt{<=}, \texttt{>}, \texttt{>=}, \texttt{in} (left-associative)
    \item \textbf{Igualdade}: \texttt{==}, \texttt{!=} (left-associative)
    \item \textbf{Bitwise AND}: \texttt{\&} (left-associative)
    \item \textbf{Bitwise XOR}: \texttt{\textasciicircum}, \texttt{\textasciicircum\textasciicircum} (left-associative)
    \item \textbf{Bitwise OR}: \texttt{|} (left-associative)
    \item \textbf{Logical AND}: \texttt{\&\&} (left-associative)
    \item \textbf{Logical OR}: \texttt{||} (left-associative)
    \item \textbf{Atribuição}: \texttt{=}, \texttt{+=}, \texttt{-=}, \texttt{*=}, \texttt{/=} (right-associative)
\end{enumerate}
\end{definition}

\subsection{Expressões Primárias}

\begin{definition}[Literais]
Expressões literais representam valores constantes:
\begin{itemize}
    \item Números: \texttt{42}, \texttt{3.14}
    \item Strings: \texttt{"text"}, \texttt{'text'}
    \item Booleanos: \texttt{true}, \texttt{false}
    \item Null: \texttt{null}
    \item Template strings: \texttt{`text \$\{expr\}`}
\end{itemize}
\end{definition}

\begin{definition}[Arrays e Tuplas Literais]
\textbf{Arrays}: \texttt{[elem1, elem2, ...]}
\begin{itemize}
    \item Podem conter elementos de tipos diferentes
    \item Suportam spread operator: \texttt{[...arr]}
\end{itemize}

\textbf{Tuplas}: \texttt{(elem1, elem2, ...)}
\begin{itemize}
    \item Tamanho fixo
    \item Tipos heterogêneos
\end{itemize}
\end{definition}

\begin{definition}[Objetos Literais]
Sintaxe: \texttt{\{ key1: value1, key2: value2, ... \}}

\textbf{Variações}:
\begin{itemize}
    \item Propriedades: \texttt{\{ x: 1, y: 2 \}}
    \item Métodos: \texttt{\{ greet() \{ ... \} \}}
    \item Chaves computadas: \texttt{\{ ["key"]: value \}}
\end{itemize}
\end{definition}

\begin{definition}[Funções Lambda (Arrow Functions)]
Sintaxe: \texttt{(params) => expression} ou \texttt{(params) => \{ statements \}}

\textbf{Formas}:
\begin{itemize}
    \item Sem parâmetros: \texttt{() => expr}
    \item Um parâmetro: \texttt{x => expr}
    \item Múltiplos: \texttt{(a, b) => expr}
    \item Com bloco: \texttt{(a, b) => \{ return a + b; \}}
    \item Rest params: \texttt{(a, ...rest) => expr}
\end{itemize}
\end{definition}

\subsection{Expressões Binárias}

\begin{definition}[Operações Aritméticas]
Semântica das operações aritméticas:
\begin{itemize}
    \item \texttt{a + b}: Se ambos números, soma. Se algum é string, concatena.
    \item \texttt{a - b}: Subtração (requer números)
    \item \texttt{a * b}: Multiplicação (requer números)
    \item \texttt{a / b}: Divisão (requer números, erro se b = 0)
    \item \texttt{a \% b}: Módulo (requer números, erro se b = 0)
    \item \texttt{a ** b}: Exponenciação (requer números)
\end{itemize}
\end{definition}

\begin{definition}[Operações Lógicas]
Operadores lógicos com short-circuit evaluation:
\begin{itemize}
    \item \texttt{a \&\& b}: Retorna \texttt{a} se falsy, senão \texttt{b}
    \item \texttt{a || b}: Retorna \texttt{a} se truthy, senão \texttt{b}
    \item \texttt{!a}: Negação booleana
\end{itemize}

\textbf{Valores Falsy}: \texttt{false}, \texttt{null}, \texttt{0}, \texttt{""}, \texttt{NaN}

\textbf{Valores Truthy}: Todos os demais
\end{definition}

\subsection{Expressões Unárias}

\begin{definition}[Operadores Unários]
\begin{itemize}
    \item \texttt{-x}: Negação numérica
    \item \texttt{+x}: Plus unário (converte para número)
    \item \texttt{!x}: NOT lógico
    \item \texttt{\textasciitilde x}: NOT bitwise
    \item \texttt{++x}: Pré-incremento
    \item \texttt{--x}: Pré-decremento
    \item \texttt{x++}: Pós-incremento
    \item \texttt{x--}: Pós-decremento
    \item \texttt{await expr}: Espera resolução de promise
\end{itemize}
\end{definition}

\subsection{Acesso a Propriedades e Chamadas}

\begin{definition}[Acesso a Membros]
\begin{itemize}
    \item \textbf{Ponto}: \texttt{obj.property}
    \item \textbf{Colchetes}: \texttt{obj["property"]}, \texttt{arr[index]}
    \item \textbf{Tupla}: \texttt{tuple.0}, \texttt{tuple.1}
    \item \textbf{Namespace}: \texttt{module::function}
\end{itemize}
\end{definition}

\begin{definition}[Chamadas de Função]
Sintaxe: \texttt{function(arg1, arg2, ...)}

\textbf{Variações}:
\begin{itemize}
    \item Chamada direta: \texttt{fn()}
    \item Chamada de método: \texttt{obj.method()}
    \item Chamada com new: \texttt{new Class()}
\end{itemize}
\end{definition}

\subsection{Match Expression}

\begin{definition}[Pattern Matching]
Sintaxe:
\begin{verbatim}
match valor {
    pattern1 => resultado1,
    pattern2 => resultado2,
    _ => default
}
\end{verbatim}

\textbf{Patterns Suportados}:
\begin{itemize}
    \item Literais: \texttt{1}, \texttt{"text"}, \texttt{true}
    \item Wildcard: \texttt{\_}
    \item Guards: condições opcionais após \texttt{=>}
\end{itemize}
\end{definition}

\section{Statements (Instruções)}

\subsection{Declarações de Variáveis}

\begin{definition}[Let Statement]
Declara variável mutável com escopo de bloco.

\textbf{Sintaxe}: \texttt{let identifier [: type] [= value];}

\textbf{Propriedades}:
\begin{itemize}
    \item Escopo de bloco (limitado ao bloco \texttt{\{\}} mais próximo)
    \item Pode ser reatribuída
    \item Se não inicializada, valor é \texttt{undefined}/\texttt{null}
    \item Múltiplas declarações: \texttt{let a = 1, b = 2;}
\end{itemize}
\end{definition}

\begin{definition}[Const Statement]
Declara variável imutável com escopo de bloco.

\textbf{Sintaxe}: \texttt{const identifier [: type] = value;}

\textbf{Propriedades}:
\begin{itemize}
    \item Escopo de bloco
    \item Deve ser inicializada na declaração
    \item Não pode ser reatribuída
    \item Para objetos/arrays, o conteúdo pode ser modificado
\end{itemize}
\end{definition}

\subsection{Statements de Controle de Fluxo}

\begin{definition}[If Statement]
Execução condicional.

\textbf{Sintaxe}:
\begin{verbatim}
if (condition) {
    // then block
} else if (condition2) {
    // else-if block
} else {
    // else block
}
\end{verbatim}

\textbf{Semântica}:
\begin{itemize}
    \item Condição é avaliada como truthy/falsy
    \item Apenas o primeiro bloco verdadeiro é executado
    \item \texttt{else} e \texttt{else if} são opcionais
\end{itemize}
\end{definition}

\begin{definition}[Loops]
\textbf{While Loop}:
\begin{verbatim}
while (condition) {
    // body
}
\end{verbatim}
Executa enquanto condição for truthy.

\textbf{Do-While Loop}:
\begin{verbatim}
do {
    // body
} while (condition);
\end{verbatim}
Executa pelo menos uma vez, depois verifica condição.

\textbf{For Loop (C-style)}:
\begin{verbatim}
for (init; condition; update) {
    // body
}
\end{verbatim}
Estrutura clássica de iteração.

\textbf{For-Each Loop}:
\begin{verbatim}
for (element in iterable) {
    // body
}
\end{verbatim}
Itera sobre elementos de coleção.
\end{definition}

\begin{definition}[Break e Continue]
\begin{itemize}
    \item \texttt{break}: Sai imediatamente do loop atual
    \item \texttt{continue}: Pula para próxima iteração do loop
\end{itemize}
Ambos afetam apenas o loop mais interno.
\end{definition}

\subsection{Tratamento de Exceções}

\begin{definition}[Try-Catch-Finally]
Sintaxe:
\begin{verbatim}
try {
    // código que pode lançar exceção
} catch (error) {
    // tratamento de erro
} finally {
    // sempre executado
}
\end{verbatim}

\textbf{Semântica}:
\begin{itemize}
    \item \texttt{try}: bloco monitorado
    \item \texttt{catch}: captura exceção, bind variável \texttt{error}
    \item \texttt{finally}: executa sempre (sucesso ou erro)
    \item Pelo menos \texttt{catch} ou \texttt{finally} deve estar presente
\end{itemize}
\end{definition}

\begin{definition}[Throw Statement]
Sintaxe: \texttt{throw expression;}

Lança uma exceção que propaga pela call stack até ser capturada por \texttt{catch}.
\end{definition}

\subsection{Return Statement}

\begin{definition}[Return]
Sintaxe: \texttt{return [expression];}

\textbf{Comportamento}:
\begin{itemize}
    \item Sai imediatamente da função
    \item Retorna valor especificado (ou \texttt{null} se omitido)
    \item Válido apenas dentro de funções ou threads
\end{itemize}
\end{definition}

\section{Declarações}

\subsection{Funções}

\begin{definition}[Function Declaration]
\textbf{Sintaxe Padrão}:
\begin{verbatim}
function name(param1 [: type1], param2 [: type2], ...)
    [: returnType] {
    // body
    return value;
}
\end{verbatim}

\textbf{Propriedades}:
\begin{itemize}
    \item Hoisted (pode ser chamada antes da declaração)
    \item Cria closure sobre escopo circundante
    \item Parâmetros podem ter valores default: \texttt{param = defaultValue}
    \item Suporta rest parameters: \texttt{...rest}
\end{itemize}
\end{definition}

\begin{definition}[Arrow Functions]
\textbf{Sintaxe}:
\begin{verbatim}
const name = (params) => expression;
const name = (params) => { statements };
\end{verbatim}

\textbf{Diferenças de \texttt{function}}:
\begin{itemize}
    \item Não são hoisted
    \item Sintaxe mais concisa
    \item Return implícito se corpo é expressão única
\end{itemize}
\end{definition}

\begin{definition}[Async Functions]
\textbf{Sintaxe}:
\begin{verbatim}
async function name(params) {
    let result = await asyncOperation();
    return result;
}
\end{verbatim}

\textbf{Comportamento}:
\begin{itemize}
    \item Pode usar \texttt{await} para esperar promises
    \item Retorna implicitamente uma promise
    \item Execução pode ser pausada em \texttt{await}
\end{itemize}
\end{definition}

\begin{definition}[Thread Functions]
\textbf{Sintaxe}:
\begin{verbatim}
thread function name(params) {
    // executa em thread separada
    return result;
}
\end{verbatim}

\textbf{Invocação}:
\begin{verbatim}
let future = thread name(args);
\end{verbatim}

\textbf{Comportamento}:
\begin{itemize}
    \item Executa em thread de OS separada
    \item Retorna handle para recuperar resultado
    \item Isolamento de contexto
\end{itemize}
\end{definition}

\subsection{Classes}

\begin{definition}[Class Declaration]
\textbf{Sintaxe Básica}:
\begin{verbatim}
class ClassName {
    // Propriedades de instância
    property = defaultValue;
    
    // Construtor
    constructor(params) {
        this.property = value;
    }
    
    // Métodos de instância
    method() {
        return this.property;
    }
    
    // Métodos estáticos
    static staticMethod() {
        return value;
    }
}
\end{verbatim}
\end{definition}

\begin{definition}[Herança]
\textbf{Sintaxe}:
\begin{verbatim}
class Child extends Parent {
    constructor(params) {
        super(parentParams);
        // inicialização adicional
    }
    
    method() {
        super.method(); // chama método do pai
        // implementação adicional
    }
}
\end{verbatim}

\textbf{Propriedades}:
\begin{itemize}
    \item Suporta herança simples (um único pai)
    \item \texttt{super} referencia a classe pai
    \item \texttt{this} referencia a instância atual
\end{itemize}
\end{definition}

\begin{definition}[Modificadores de Acesso]
\begin{itemize}
    \item \textbf{public}: Acessível de qualquer lugar (default)
    \item \textbf{private}: Acessível apenas dentro da classe
    \item \textbf{protected}: Acessível na classe e subclasses
\end{itemize}

\textbf{Aplicação}:
\begin{verbatim}
class Example {
    public x = 1;      // acessível externamente
    private y = 2;     // apenas interno
    protected z = 3;   // classe e filhas
}
\end{verbatim}
\end{definition}

\begin{definition}[Membros Estáticos]
Membros declarados com \texttt{static} pertencem à classe, não às instâncias.

\textbf{Exemplo}:
\begin{verbatim}
class Counter {
    static count = 0;
    
    static increment() {
        Counter.count++;
    }
}

// Uso
Counter.increment();
console.log(Counter.count); // 1
\end{verbatim}
\end{definition}

\section{Sistema de Módulos}

\subsection{Import Statement}

\begin{definition}[Importações Nomeadas]
Sintaxe: \texttt{import \{ name1, name2, ... \} from "module";}

Importa exports nomeados de um módulo.
\end{definition}

\begin{definition}[Importação Namespace]
Sintaxe: \texttt{import * as namespace from "module";}

Importa todos os exports sob um namespace.

\textbf{Uso}: \texttt{namespace.function()}
\end{definition}

\begin{definition}[Importação Side-Effect]
Sintaxe: \texttt{import "module";}

Executa o módulo por efeitos colaterais, sem importar valores.
\end{definition}

\begin{definition}[Importação por Identificador]
Sintaxe: \texttt{import moduleName;}

Importa módulo sem quotes (syntax moderna).

\textbf{Acesso com namespace operator}:
\begin{verbatim}
import ipe;
ipe::init();
ipe::config::value;
\end{verbatim}
\end{definition}

\subsection{Export Statement}

\begin{definition}[Exports]
\textbf{Export Direto}:
\begin{verbatim}
export function name() { }
export const VALUE = 42;
export class ClassName { }
\end{verbatim}

\textbf{Re-export}:
\begin{verbatim}
export { name1, name2 } from "other_module";
export * from "other_module";
\end{verbatim}
\end{definition}

\subsection{Native Directives}

\begin{definition}[Carregamento de Módulos Nativos]
Sintaxe: \texttt{\#<module\_name>}

Carrega módulo nativo (implementado em C/Rust) no escopo global.

\textbf{Módulos Disponíveis}:
\begin{itemize}[noitemsep]
    \item \texttt{\#io} - Operações de I/O
    \item \texttt{\#crypto} - Criptografia
    \item \texttt{\#http} - Cliente/servidor HTTP
    \item \texttt{\#math} - Funções matemáticas
    \item \texttt{\#tcp}, \texttt{\#udp} - Networking
    \item \texttt{\#time} - Data e hora
    \item Outros módulos nativos
\end{itemize}

\textbf{Exemplo}:
\begin{verbatim}
#io

let content = io_read_file("data.txt");
io_write_file("output.txt", content);
\end{verbatim}
\end{definition}

\section{Modelo de Execução}

\subsection{Arquitetura de Execução Híbrida}

\begin{definition}[Três Modos de Execução]
Dryad suporta três estratégias de execução:

\begin{enumerate}
    \item \textbf{Tree-Walking Interpreter}:
    \begin{itemize}
        \item Executa AST diretamente
        \item Modo padrão
        \item Otimizado para desenvolvimento e debugging
        \item Mais lento, mas mais flexível
    \end{itemize}
    
    \item \textbf{Bytecode VM}:
    \begin{itemize}
        \item Compila AST para bytecode stack-based
        \item Habilitado com \texttt{set\_compile\_mode(true)}
        \item VM executa OpCodes sequencialmente
        \item Performance intermediária
    \end{itemize}
    
    \item \textbf{AOT Compiler} (Experimental):
    \begin{itemize}
        \item Compila bytecode/IR para código nativo
        \item Gera executáveis nativos (ARM64, x86\_64)
        \item Suporta formatos ELF e PE
        \item Máxima performance
    \end{itemize}
\end{enumerate}
\end{definition}

\subsection{Pipeline de Compilação}

\begin{definition}[Fluxo de Compilação]
O processo de compilação segue:

\[
\text{Source Code} \xrightarrow{\text{Lexer}} \text{Tokens} \xrightarrow{\text{Parser}} \text{AST}
\]

\[
\text{AST} \xrightarrow{\text{Interpreter}} \text{Result}
\]

ou

\[
\text{AST} \xrightarrow{\text{Bytecode Compiler}} \text{OpCodes} \xrightarrow{\text{VM}} \text{Result}
\]

ou

\[
\text{OpCodes} \xrightarrow{\text{IR Converter}} \text{IR} \xrightarrow{\text{Backend}} \text{Assembly} \xrightarrow{\text{Generator}} \text{Native Binary}
\]
\end{definition}

\subsection{Escopo e Resolução de Variáveis}

\begin{definition}[Cadeia de Escopos]
A hierarquia de escopos em Dryad:

\begin{verbatim}
Global Scope
  ├─ Module Scope
  │   ├─ Function Scope
  │   │   ├─ Block Scope
  │   │   ├─ Block Scope
  │   │   └─ ...
  │   └─ ...
  ├─ Class Scope
  │   └─ Method Scope
  └─ ...
\end{verbatim}

\textbf{Algoritmo de Lookup}:
\begin{enumerate}
    \item Verifica escopo atual
    \item Se não encontrado, verifica escopo pai
    \item Continua até escopo global
    \item Se não encontrado em nenhum escopo, lança \texttt{ReferenceError}
\end{enumerate}
\end{definition}

\begin{property}[Regras de Binding]
\begin{itemize}
    \item \texttt{let} e \texttt{const}: block-scoped
    \item \texttt{function}: hoisted para escopo de função/módulo
    \item Variáveis globais: acessíveis de qualquer lugar
    \item Closures: funções capturam variáveis do escopo de definição
\end{itemize}
\end{property}

\subsection{Execução de Funções}

\begin{definition}[Chamada de Função]
Ao chamar uma função:

\begin{enumerate}
    \item Cria novo escopo de função
    \item Faz binding de parâmetros com argumentos
    \item Inicializa variáveis declaradas no corpo
    \item Executa statements em ordem
    \item Retorna valor (explícito via \texttt{return} ou \texttt{null} implícito)
    \item Destrói escopo de função
\end{enumerate}
\end{definition}

\begin{definition}[Closures]
Funções capturam variáveis do escopo de definição.

\textbf{Propriedades}:
\begin{itemize}
    \item Variáveis capturadas permanecem vivas enquanto função existir
    \item Modificações em variáveis capturadas são visíveis
    \item Múltiplas closures podem compartilhar mesmas variáveis
\end{itemize}

\textbf{Exemplo Conceitual}:
\begin{verbatim}
function makeCounter() {
    let count = 0;  // capturada por closure
    return function() {
        return ++count;  // modifica variável capturada
    };
}
\end{verbatim}
\end{definition}

\subsection{Instanciação de Classes}

\begin{definition}[Processo new ClassName()]
Ao executar \texttt{new ClassName(args)}:

\begin{enumerate}
    \item Cria nova instância de objeto
    \item Se \texttt{constructor} existe:
    \begin{enumerate}
        \item Chama constructor com \texttt{this} bound à instância
        \item Executa corpo do constructor
    \end{enumerate}
    \item Retorna a instância
\end{enumerate}
\end{definition}

\begin{definition}[Execução de Métodos]
Ao chamar \texttt{obj.method(args)}:

\begin{enumerate}
    \item Cria escopo de função
    \item Faz binding de \texttt{this} ao objeto
    \item Faz binding de parâmetros com argumentos
    \item Executa corpo do método
    \item Retorna resultado
\end{enumerate}
\end{definition}

\subsection{Representação Interna de Valores}

\begin{definition}[Value Enum (Conceitual)]
Valores em runtime são representados por variante enum:

\begin{itemize}
    \item \texttt{Number(f64)} - IEEE 754 floating-point
    \item \texttt{String(String)} - UTF-8 string
    \item \texttt{Bool(bool)} - booleano
    \item \texttt{Null} - valor nulo
    \item \texttt{Array(Vec<Value>)} - array dinâmico
    \item \texttt{Object(HashMap<String, Value>)} - mapa chave-valor
    \item \texttt{Function(...)} - referência a função
    \item \texttt{Class(...)} - definição ou instância de classe
    \item \texttt{Thread(...)} - handle de thread
\end{itemize}
\end{definition}

\subsection{Gerenciamento de Memória}

\begin{property}[Garbage Collection]
\begin{itemize}
    \item \textbf{Estratégia}: Reference counting (via Rust)
    \item \textbf{Automático}: Sem intervenção do programador
    \item \textbf{Stack}: Valores primitivos permanecem na pilha
    \item \textbf{Heap}: Tipos complexos alocados no heap
    \item \textbf{Cleanup}: Valores liberados quando sem referências
    \item \textbf{Transparência}: Usuário não gerencia memória manualmente
\end{itemize}
\end{property}

\section{Sistema de Erros}

\subsection{Categorias de Erros}

\begin{definition}[DryadError]
Todos os erros em Dryad são instâncias de \texttt{DryadError} com:

\begin{itemize}
    \item \textbf{Código}: Identificador numérico único
    \item \textbf{Categoria}: Tipo de erro
    \item \textbf{Mensagem}: Descrição legível
    \item \textbf{Localização}: Arquivo, linha, coluna
    \item \textbf{Stack Trace}: Pilha de chamadas
    \item \textbf{Contexto de Debug}: Sugestões, variáveis locais, documentação
\end{itemize}
\end{definition}

\begin{definition}[Categorias de Erro]
\begin{enumerate}
    \item \textbf{Lexer} (1000-1999):
    \begin{itemize}[noitemsep]
        \item Caractere inesperado
        \item String não fechada
        \item Número inválido
        \item Escape sequence inválida
    \end{itemize}
    
    \item \textbf{Parser} (2000-2999):
    \begin{itemize}[noitemsep]
        \item Token inesperado
        \item Estrutura sintática inválida
        \item Parênteses não balanceados
        \item Statement esperado
    \end{itemize}
    
    \item \textbf{Runtime} (3000-3999):
    \begin{itemize}[noitemsep]
        \item Variável indefinida
        \item Tipo incompatível
        \item Divisão por zero
        \item Stack overflow
        \item Propriedade não encontrada
    \end{itemize}
    
    \item \textbf{Type} (4000-4999):
    \begin{itemize}[noitemsep]
        \item Incompatibilidade de tipos
        \item Número incorreto de argumentos
    \end{itemize}
    
    \item \textbf{I/O} (5000-5999):
    \begin{itemize}[noitemsep]
        \item Arquivo não encontrado
        \item Permissão negada
        \item Operação de I/O falhou
    \end{itemize}
    
    \item \textbf{Module} (6000-6999):
    \begin{itemize}[noitemsep]
        \item Módulo não encontrado
        \item Módulo nativo desconhecido
        \item Erro de resolução
    \end{itemize}
    
    \item \textbf{Syntax} (7000-7999):
    \begin{itemize}[noitemsep]
        \item Erro de sintaxe geral
    \end{itemize}
    
    \item \textbf{Warning} (8000-8999):
    \begin{itemize}[noitemsep]
        \item Avisos não-fatais
        \item Severidade: Low, Medium, High
    \end{itemize}
    
    \item \textbf{System} (9000+):
    \begin{itemize}[noitemsep]
        \item Erros de sistema
    \end{itemize}
\end{enumerate}
\end{definition}

\subsection{Propagação de Exceções}

\begin{property}[Propagação]
Exceções não capturadas propagam pela call stack:

\begin{enumerate}
    \item Exceção lançada em função profunda
    \item Se não há \texttt{try-catch} local, propaga para chamador
    \item Continua até encontrar \texttt{catch} ou chegar ao topo
    \item Se não capturada, termina programa com erro
\end{enumerate}
\end{property}

\begin{definition}[Stack Trace]
Cada exceção carrega stack trace com:
\begin{itemize}
    \item Nome da função
    \item Localização (arquivo:linha:coluna)
    \item Contexto adicional
\end{itemize}

Ordenado do mais profundo ao mais superficial.
\end{definition}

\section{Modelo de Concorrência}

\subsection{Async/Await}

\begin{definition}[Funções Assíncronas]
Funções declaradas com \texttt{async} podem usar \texttt{await}.

\textbf{Comportamento}:
\begin{itemize}
    \item \texttt{async function} retorna promise implicitamente
    \item \texttt{await expr} suspende execução até promise resolver
    \item Permite código assíncrono com sintaxe síncrona
\end{itemize}

\textbf{Modelo Conceitual}:
\begin{verbatim}
async function process() {
    let data = await fetchData();  // pausa aqui
    return processData(data);      // continua quando resolve
}
\end{verbatim}
\end{definition}

\subsection{Threads Nativos}

\begin{definition}[Thread Functions]
Funções declaradas com \texttt{thread function} podem executar em threads de OS.

\textbf{Propriedades}:
\begin{itemize}
    \item Executa em thread separada do SO
    \item Isolamento de contexto
    \item Comunicação via mutexes ou retorno
    \item Handle retornado permite recuperar resultado
\end{itemize}

\textbf{Exemplo Conceitual}:
\begin{verbatim}
thread function backgroundTask(id) {
    // executa em paralelo
    return result;
}

let handle1 = thread backgroundTask(1);
let handle2 = thread backgroundTask(2);
// ambos executam simultaneamente
\end{verbatim}
\end{definition}

\subsection{Sincronização com Mutex}

\begin{definition}[Mutex]
Primitiva de sincronização para proteger recursos compartilhados.

\textbf{API}:
\begin{itemize}
    \item \texttt{mutex()}: cria novo mutex
    \item \texttt{mu.lock()}: adquire lock (bloqueia se ocupado)
    \item \texttt{mu.unlock()}: libera lock
\end{itemize}

\textbf{Uso}:
\begin{verbatim}
let mu = mutex();
let sharedData = 0;

thread function increment() {
    mu.lock();
    sharedData++;  // seção crítica
    mu.unlock();
}
\end{verbatim}
\end{definition}

\section{Foreign Function Interface (FFI)}

\subsection{Sistema de Módulos Nativos}

\begin{definition}[Native Modules]
Módulos implementados em C/Rust e expostos para Dryad.

\textbf{Características}:
\begin{itemize}
    \item Carregados via native directives (\texttt{\#<module>})
    \item Funções tornam-se disponíveis no escopo global
    \item Tipagem dinâmica na interface
    \item Convenção de chamada C
\end{itemize}
\end{definition}

\subsection{Mapeamento de Tipos}

\begin{definition}[Type Mapping Dryad ↔ C]
\begin{center}
\begin{tabular}{|l|l|l|}
\hline
\textbf{Dryad Type} & \textbf{C Type} & \textbf{Notas} \\
\hline
\texttt{number} & \texttt{double} (f64) & IEEE 754 \\
\texttt{string} & \texttt{const char*} & UTF-8 null-terminated \\
\texttt{bool} & \texttt{bool} / \texttt{uint8\_t} & Booleano C \\
\texttt{null} & \texttt{NULL} / \texttt{nullptr} & Ponteiro nulo \\
\texttt{array} & \texttt{void*} & Ponteiro opaco \\
\texttt{object} & \texttt{void*} & Ponteiro opaco \\
\hline
\end{tabular}
\end{center}

\textbf{Limitações}:
\begin{itemize}
    \item Tipos complexos passados como ponteiros opacos
    \item Apenas primitivos têm conversão direta
    \item Responsabilidade de conversão é da implementação nativa
\end{itemize}
\end{definition}

\subsection{Convenção de Chamada}

\begin{property}[FFI Call Convention]
Chamadas FFI seguem convenção C:
\begin{itemize}
    \item Parâmetros passados left-to-right
    \item Stack-based (x86/x64)
    \item Valores de retorno em registrador \texttt{rax} (x86-64)
    \item Caller limpa stack (cdecl)
\end{itemize}
\end{property}

\section{Arquitetura de Runtime Mínimo com Intrinsics}

\subsection{Filosofia do Self-Hosting}

\begin{definition}[Runtime Mínimo]
Dryad adota uma filosofia de \textbf{self-hosting} onde:
\begin{itemize}
    \item Runtime C++ fornece apenas \textbf{syscalls primitivas} ($\sim$50 intrinsics)
    \item Standard library implementada \textbf{100\% em Dryad}
    \item Zero necessidade de escrever bindings manualmente
    \item Bibliotecas de alto nível (HTTP, JSON, crypto) escritas em Dryad puro
\end{itemize}
\end{definition}

\begin{property}[Vantagens Arquiteturais]
Esta abordagem, utilizada por linguagens modernas como Go, Zig e Rust, oferece:
\begin{itemize}
    \item \textbf{Manutenção Única}: Não duplicar código em C++ e Dryad
    \item \textbf{Extensibilidade}: Usuários implementam bibliotecas em Dryad
    \item \textbf{Testabilidade}: Mock implementations sem tocar código nativo
    \item \textbf{Debuggability}: Todo código visível e inspecionável
    \item \textbf{Portabilidade}: Apenas syscalls precisam ser adaptadas por OS
\end{itemize}
\end{property}

\subsection{Sistema de Intrinsics}

\begin{definition}[Intrinsic Function]
Uma \textbf{intrinsic} é uma função especial reconhecida pelo compilador que:
\begin{enumerate}
    \item Não gera instrução normal de chamada (\texttt{CALL})
    \item Compila diretamente para instrução assembly/opcode especializado
    \item Tem custo de execução próximo a zero
    \item Mapeia diretamente para código C++ do runtime
\end{enumerate}
\end{definition}

\begin{definition}[Sintaxe de Declaração]
Intrinsics são declaradas com decorator \texttt{@intrinsic}:

\begin{verbatim}
// Declaração de syscall intrinsic
@intrinsic("syscall.open")
extern function __open(path: string, flags: i32): i32;

@intrinsic("syscall.read")
extern function __read(fd: i32, buf: ptr<u8>, len: usize): isize;

@intrinsic("syscall.write")
extern function __write(fd: i32, buf: ptr<u8>, len: usize): isize;

@intrinsic("syscall.close")
extern function __close(fd: i32): void;
\end{verbatim}
\end{definition}

\subsection{Pipeline de Compilação com Intrinsics}

\begin{definition}[Processo de Compilação]
Durante compilação, ao encontrar chamada de função com \texttt{@intrinsic}:

\begin{enumerate}
    \item \textbf{Parser}: Reconhece decorator e marca função como intrinsic
    \item \textbf{Type Checker}: Valida assinatura contra definição do runtime
    \item \textbf{Code Generator}: Emite opcode especial ao invés de \texttt{CALL}
    \item \textbf{VM/AOT}: Executa código C++ diretamente
\end{enumerate}

\textbf{Exemplo de bytecode gerado}:
\begin{verbatim}
// Dryad source
let fd = __open("file.txt", O_RDONLY);

// Bytecode emitido
LOAD_CONST "file.txt"
LOAD_CONST 0              // O_RDONLY
INTRINSIC_SYSCALL 1       // ID 1 = syscall.open
STORE_LOCAL fd
\end{verbatim}
\end{definition}

\subsection{Implementação no Runtime C++}

\begin{definition}[Tabela de Syscalls]
Runtime mantém tabela de IDs de syscalls:

\begin{verbatim}
// C++ implementation
enum class SyscallID : uint16_t {
    OPEN = 1,
    READ = 2,
    WRITE = 3,
    CLOSE = 4,
    SOCKET = 5,
    CONNECT = 6,
    // ... ~50 syscalls total
};
\end{verbatim}
\end{definition}

\begin{definition}[Dispatcher de Intrinsics]
VM executa intrinsics via dispatch table:

\begin{verbatim}
void VM::execute_intrinsic(SyscallID id) {
    switch (id) {
        case SyscallID::OPEN: {
            auto flags = pop_stack().as_int();
            auto path = pop_stack().as_string();
            int fd = ::open(path.c_str(), flags);
            push_stack(Value::from_int(fd));
            break;
        }
        case SyscallID::READ: {
            auto len = pop_stack().as_int();
            auto buf = pop_stack().as_buffer();
            auto fd = pop_stack().as_int();
            ssize_t n = ::read(fd, buf.data(), len);
            push_stack(Value::from_int(n));
            break;
        }
        // ... outros syscalls
    }
}
\end{verbatim}
\end{definition}

\subsection{Standard Library em Dryad Puro}

\begin{definition}[Camadas de Abstração]
Standard library organizada em camadas:

\begin{verbatim}
┌────────────────────────────────────────┐
│  High-Level APIs (100% Dryad)          │
│  - http.dryad (HTTP client/server)     │
│  - json.dryad (parser/serializer)      │
│  - crypto.dryad (algorithms)           │
└────────────────────────────────────────┘
              ↓ usa
┌────────────────────────────────────────┐
│  Core I/O (Dryad + Intrinsics)         │
│  - io.dryad (FileSystem, Buffer)       │
│  - net.dryad (Socket, TCP/UDP)         │
└────────────────────────────────────────┘
              ↓ chama
┌────────────────────────────────────────┐
│  Runtime Intrinsics (~50 syscalls)     │
│  - syscall.open/read/write/close       │
│  - syscall.socket/connect/send/recv    │
│  - memory.alloc/free                   │
└────────────────────────────────────────┘
\end{verbatim}
\end{definition}

\begin{definition}[Exemplo: FileSystem em Dryad]
Implementação completa de I/O em Dryad puro:

\begin{verbatim}
// @std/io.dryad
@intrinsic("syscall.open")
extern function __open(path: string, flags: i32): i32;

@intrinsic("syscall.read")
extern function __read(fd: i32, buf: ptr<u8>, len: usize): isize;

@intrinsic("syscall.close")
extern function __close(fd: i32): void;

const O_RDONLY = 0;
const O_WRONLY = 1;
const O_CREAT = 64;

export class File {
    private fd: number;
    
    constructor(path: string, mode: string = "r") {
        let flags = mode == "r" ? O_RDONLY : (O_WRONLY | O_CREAT);
        this.fd = __open(path, flags);
        if (this.fd < 0) {
            throw new Error("Failed to open: " + path);
        }
    }
    
    read(): string {
        let buffer = Buffer.allocate(4096);
        let n = __read(this.fd, buffer.ptr, 4096);
        if (n < 0) throw new Error("Read failed");
        return buffer.slice(0, n).toString();
    }
    
    close(): void {
        __close(this.fd);
    }
}

export function readFile(path: string): string {
    let f = new File(path, "r");
    let content = f.read();
    f.close();
    return content;
}
\end{verbatim}
\end{definition}

\subsection{Virtual File System (VFS)}

\begin{definition}[Arquitetura VFS]
Sistema de arquivos virtual com backends plugáveis:

\begin{verbatim}
interface FileSystemBackend {
    open(path: string, mode: string): FileHandle;
    read(handle: FileHandle, size: number): Buffer;
    write(handle: FileHandle, data: Buffer): number;
    close(handle: FileHandle): void;
}

// Backend nativo (usa syscalls)
class NativeBackend implements FileSystemBackend {
    open(path: string, mode: string): FileHandle {
        let flags = this.parseMode(mode);
        let fd = __open(path, flags);
        return new NativeFileHandle(fd);
    }
    // ... implementação com syscalls
}

// Backend in-memory (100% Dryad, zero syscalls!)
class MemoryBackend implements FileSystemBackend {
    private storage: Map<string, Buffer> = new Map();
    
    open(path: string, mode: string): FileHandle {
        if (!this.storage.has(path) && mode == "w") {
            this.storage.set(path, Buffer.allocate(0));
        }
        return new MemoryFileHandle(path, this.storage);
    }
    // ... sem syscalls, apenas Map em memória
}
\end{verbatim}
\end{definition}

\begin{property}[Benefícios do VFS]
\begin{itemize}
    \item \textbf{Testabilidade}: Testes usam \texttt{MemoryBackend} sem tocar disco
    \item \textbf{Extensibilidade}: Usuários criam backends customizados (HTTP, S3, etc)
    \item \textbf{Segurança}: Backends podem implementar sandboxing
    \item \textbf{Composabilidade}: Backends podem ser combinados (cache, logging)
\end{itemize}
\end{property}

\subsection{Networking e Async I/O}

\begin{definition}[Event Loop em Dryad]
Para suportar operações assíncronas, event loop implementado em Dryad:

\begin{verbatim}
// @std/async/event_loop.dryad
@intrinsic("syscall.epoll_create")
extern function __epoll_create(): i32;

@intrinsic("syscall.epoll_wait")
extern function __epoll_wait(epfd: i32, timeout: i32): Array<i32>;

@intrinsic("syscall.epoll_ctl")
extern function __epoll_ctl(epfd: i32, op: i32, fd: i32): void;

class EventLoop {
    private epoll_fd: number;
    private tasks: Map<number, Task> = new Map();
    
    constructor() {
        this.epoll_fd = __epoll_create();
    }
    
    run(): void {
        while (this.tasks.size > 0) {
            // Multiplexação de I/O não-bloqueante
            let ready_fds = __epoll_wait(this.epoll_fd, 100);
            
            for (let fd of ready_fds) {
                let task = this.tasks.get(fd);
                if (task) {
                    task.resume(); // Acorda coroutine
                }
            }
        }
    }
    
    register(fd: number, task: Task): void {
        __epoll_ctl(this.epoll_fd, EPOLL_CTL_ADD, fd);
        this.tasks.set(fd, task);
    }
}
\end{verbatim}
\end{definition}

\begin{definition}[HTTP Client em Dryad Puro]
Cliente HTTP implementado sem código C++:

\begin{verbatim}
// @std/http.dryad
import { Socket } from "@std/net";
import { EventLoop } from "@std/async";

export class HttpClient {
    async get(url: string): Response {
        // Parse URL (puro Dryad)
        let parsed = this.parseUrl(url);
        
        // Socket connection (syscall socket/connect)
        let socket = new Socket(parsed.host, parsed.port);
        await socket.connect();
        
        // Build HTTP request (puro Dryad)
        let request = this.buildRequest("GET", parsed.path);
        
        // Send (syscall write)
        await socket.write(request);
        
        // Receive (syscall read via event loop)
        let response = await socket.readAll();
        
        // Parse response (puro Dryad)
        return this.parseResponse(response);
    }
    
    private parseUrl(url: string): UrlInfo { ... }
    private buildRequest(method: string, path: string): string { ... }
    private parseResponse(data: string): Response { ... }
}
\end{verbatim}
\end{definition}

\subsection{Gerenciamento de Memória com Buffer}

\begin{definition}[Classe Buffer]
Abstração type-safe para manipulação de bytes:

\begin{verbatim}
// @std/buffer.dryad
@intrinsic("memory.allocate")
extern function __alloc(size: usize): ptr<u8>;

@intrinsic("memory.free")
extern function __free(ptr: ptr<u8>): void;

export class Buffer {
    private data: ptr<u8>;
    private _size: number;
    
    static allocate(size: number): Buffer {
        let ptr = __alloc(size);
        return new Buffer(ptr, size);
    }
    
    static fromString(s: string): Buffer {
        let size = s.length;
        let buf = Buffer.allocate(size);
        // Copia bytes de string para buffer
        for (let i = 0; i < size; i++) {
            buf.set(i, s.charCodeAt(i));
        }
        return buf;
    }
    
    toString(): string {
        // Converte bytes para string
        let chars = [];
        for (let i = 0; i < this._size; i++) {
            chars.push(String.fromCharCode(this.get(i)));
        }
        return chars.join("");
    }
    
    get ptr(): ptr<u8> { return this.data; }
    get size(): number { return this._size; }
}
\end{verbatim}
\end{definition}

\begin{property}[Segurança de Memória]
Buffer encapsula ponteiros brutos e fornece:
\begin{itemize}
    \item \textbf{Bounds checking}: Verificação de limites em runtime
    \item \textbf{Automatic cleanup}: GC libera memória automaticamente
    \item \textbf{Type safety}: Conversões explícitas entre Buffer e tipos primitivos
    \item \textbf{Zero-copy}: Passa ponteiro direto para syscalls
\end{itemize}
\end{property}

\subsection{Conjunto Mínimo de Syscalls}

\begin{definition}[Syscalls Essenciais]
Runtime C++ implementa aproximadamente 50 syscalls:

\textbf{File I/O}:
\begin{itemize}[noitemsep]
    \item \texttt{open}, \texttt{read}, \texttt{write}, \texttt{close}
    \item \texttt{lseek}, \texttt{stat}, \texttt{unlink}, \texttt{mkdir}
\end{itemize}

\textbf{Network}:
\begin{itemize}[noitemsep]
    \item \texttt{socket}, \texttt{connect}, \texttt{bind}, \texttt{listen}
    \item \texttt{accept}, \texttt{send}, \texttt{recv}, \texttt{shutdown}
\end{itemize}

\textbf{Memory}:
\begin{itemize}[noitemsep]
    \item \texttt{malloc}, \texttt{free}, \texttt{realloc}
    \item \texttt{memcpy}, \texttt{memset}
\end{itemize}

\textbf{Async I/O}:
\begin{itemize}[noitemsep]
    \item \texttt{epoll\_create}, \texttt{epoll\_ctl}, \texttt{epoll\_wait} (Linux)
    \item \texttt{kqueue}, \texttt{kevent} (macOS/BSD)
    \item \texttt{select} (fallback portável)
\end{itemize}

\textbf{Process/Thread}:
\begin{itemize}[noitemsep]
    \item \texttt{fork}, \texttt{exec}, \texttt{wait}
    \item \texttt{pthread\_create}, \texttt{pthread\_join}
\end{itemize}

\textbf{Time}:
\begin{itemize}[noitemsep]
    \item \texttt{gettimeofday}, \texttt{clock\_gettime}
    \item \texttt{sleep}, \texttt{nanosleep}
\end{itemize}
\end{definition}

\subsection{Comparação com Sistema Anterior}

\begin{center}
\begin{tabular}{|l|p{5cm}|p{5cm}|}
\hline
\textbf{Aspecto} & \textbf{Sistema Antigo} & \textbf{Sistema Novo} \\
\hline
Stdlib & Wrappers C++ manuais & 100\% Dryad \\
\hline
Manutenção & Duplicada (C++ + Dryad) & Única (apenas Dryad) \\
\hline
Extensibilidade & Requer código C++ & Bibliotecas em Dryad \\
\hline
Testabilidade & Difícil (syscalls reais) & Backends mock \\
\hline
Runtime Size & Grande (muitos bindings) & Mínimo (~50 syscalls) \\
\hline
Performance & Overhead de wrapper & Intrinsics diretas \\
\hline
Portabilidade & Adaptar cada binding & Adaptar syscalls \\
\hline
Debuggability & Parte em C++ opaco & Tudo visível em Dryad \\
\hline
\end{tabular}
\end{center}

\section{Bytecode e IR}

\subsection{Bytecode VM}

\begin{definition}[OpCodes]
O bytecode de Dryad consiste em OpCodes stack-based.

\textbf{Categorias de OpCodes}:
\begin{itemize}
    \item \textbf{Stack}: Push, Pop, Dup
    \item \textbf{Aritméticos}: Add, Sub, Mul, Div, Mod, Pow
    \item \textbf{Lógicos}: And, Or, Not
    \item \textbf{Bitwise}: BitwiseAnd, BitwiseOr, BitwiseXor, Shift
    \item \textbf{Comparação}: Equal, NotEqual, Less, Greater, LessEqual, GreaterEqual
    \item \textbf{Controle}: Jump, JumpIfFalse, JumpIfTrue, Call, Return
    \item \textbf{Variáveis}: GetLocal, SetLocal, GetGlobal, SetGlobal
    \item \textbf{Funções}: DefineFunction, CallFunction
    \item \textbf{Objetos}: GetProperty, SetProperty, NewObject, NewArray
    \item \textbf{Classes}: DefineClass, NewInstance, CallMethod
\end{itemize}
\end{definition}

\begin{definition}[Stack-Based Execution]
A VM mantém:
\begin{itemize}
    \item \textbf{Stack}: Pilha de valores para operandos e resultados
    \item \textbf{Instruction Pointer (IP)}: Posição atual no bytecode
    \item \textbf{Call Stack}: Pilha de frames de função
    \item \textbf{Globals}: Ambiente de variáveis globais
\end{itemize}

\textbf{Ciclo de Execução}:
\begin{enumerate}
    \item Fetch: Lê OpCode no IP
    \item Decode: Identifica operação e operandos
    \item Execute: Realiza operação (manipula stack)
    \item Advance: Incrementa IP
    \item Repeat até \texttt{Return} ou fim
\end{enumerate}
\end{definition}

\subsection{Intermediate Representation (IR)}

\begin{definition}[IR Dryad]
Representação intermediária usada pelo compilador AOT.

\textbf{Características}:
\begin{itemize}
    \item Mais abstrata que assembly, mais concreta que AST
    \item Register-based (ao contrário do bytecode stack-based)
    \item Independente de arquitetura
    \item Permite otimizações antes de geração de código
\end{itemize}

\textbf{Estrutura}:
\begin{itemize}
    \item \texttt{IrModule}: Unidade de compilação completa
    \item \texttt{IrFunction}: Definição de função com parâmetros e corpo
    \item \texttt{IrInstruction}: Instruções IR (Load, Store, Add, Call, etc.)
    \item \texttt{IrValue}: Valores (constantes, variáveis, registradores)
\end{itemize}
\end{definition}

\section{Compilação AOT (Ahead-of-Time)}

\subsection{Pipeline AOT}

\begin{definition}[Fases do Compilador AOT]
\begin{enumerate}
    \item \textbf{Bytecode → IR Conversion}:
    \begin{itemize}
        \item Converte OpCodes stack-based para IR register-based
        \item Análise de fluxo de controle
        \item Construção de CFG (Control Flow Graph)
    \end{itemize}
    
    \item \textbf{IR Optimization}:
    \begin{itemize}
        \item Constant folding
        \item Dead code elimination
        \item Register allocation
    \end{itemize}
    
    \item \textbf{Backend Code Generation}:
    \begin{itemize}
        \item Seleção de instruções para arquitetura alvo
        \item Alocação de registradores físicos
        \item Geração de assembly (ARM64 ou x86\_64)
    \end{itemize}
    
    \item \textbf{Object File Generation}:
    \begin{itemize}
        \item Geração de formato executável (ELF ou PE)
        \item Linking com runtime libraries
        \item Produção de binário nativo
    \end{itemize}
\end{enumerate}
\end{definition}

\subsection{Backends Suportados}

\begin{definition}[Arquiteturas]
\begin{itemize}
    \item \textbf{ARM64}:
    \begin{itemize}
        \item Arquitetura RISC de 64 bits
        \item Registradores: x0-x30, sp, pc
        \item Instrução fixa de 32 bits
    \end{itemize}
    
    \item \textbf{x86\_64}:
    \begin{itemize}
        \item Arquitetura CISC de 64 bits
        \item Registradores: rax, rbx, rcx, rdx, rsi, rdi, rsp, rbp, r8-r15
        \item Comprimento de instrução variável
    \end{itemize}
\end{itemize}
\end{definition}

\subsection{Formatos de Objeto}

\begin{definition}[Geradores de Formato]
\begin{itemize}
    \item \textbf{ELF Generator}:
    \begin{itemize}
        \item Executable and Linkable Format
        \item Usado em Linux/Unix
        \item Seções: .text, .data, .bss, .rodata
    \end{itemize}
    
    \item \textbf{PE Generator}:
    \begin{itemize}
        \item Portable Executable
        \item Usado em Windows
        \item Seções: .text, .data, .rdata
    \end{itemize}
\end{itemize}
\end{definition}

\section{Padrões de Programação}

\subsection{Paradigma Orientado a Objetos}

\begin{definition}[OOP em Dryad]
Dryad suporta OOP completo:
\begin{itemize}
    \item \textbf{Encapsulamento}: Modificadores de acesso
    \item \textbf{Herança}: Via \texttt{extends}
    \item \textbf{Polimorfismo}: Métodos podem ser sobrescritos
    \item \textbf{Abstração}: Classes e interfaces
\end{itemize}
\end{definition}

\subsection{Paradigma Funcional}

\begin{definition}[Functional Programming em Dryad]
Características funcionais:
\begin{itemize}
    \item \textbf{First-class functions}: Funções como valores
    \item \textbf{Higher-order functions}: Funções que recebem/retornam funções
    \item \textbf{Closures}: Captura de escopo
    \item \textbf{Arrow functions}: Sintaxe concisa
    \item \textbf{Imutabilidade parcial}: \texttt{const} para variáveis
\end{itemize}
\end{definition}

\subsection{Paradigma Procedural}

\begin{definition}[Procedural em Dryad]
Suporte a programação procedural tradicional:
\begin{itemize}
    \item Funções top-level
    \item Variáveis globais e locais
    \item Controle de fluxo estruturado
    \item Sem necessidade de classes
\end{itemize}
\end{definition}

\section{Especificações Formais}

\subsection{Gramática Formal (BNF Simplificada)}

\begin{definition}[Gramática de Dryad (Subset)]
\begin{verbatim}
<program> ::= <statement>*

<statement> ::= <var-decl>
              | <const-decl>
              | <function-decl>
              | <class-decl>
              | <if-stmt>
              | <while-stmt>
              | <for-stmt>
              | <return-stmt>
              | <throw-stmt>
              | <try-catch-stmt>
              | <expression-stmt>
              | <import-stmt>
              | <export-stmt>
              | <block-stmt>

<var-decl> ::= "let" <identifier> [":" <type>] ["=" <expression>] ";"

<const-decl> ::= "const" <identifier> [":" <type>] "=" <expression> ";"

<function-decl> ::= ["async"] "function" <identifier>
                    "(" <parameter-list> ")" [":" <type>]
                    "{" <statement>* "}"

<class-decl> ::= "class" <identifier> ["extends" <identifier>]
                 "{" <class-member>* "}"

<expression> ::= <literal>
               | <identifier>
               | <binary-expr>
               | <unary-expr>
               | <call-expr>
               | <member-expr>
               | <lambda-expr>
               | <array-literal>
               | <object-literal>
               | <match-expr>

<binary-expr> ::= <expression> <binary-op> <expression>

<unary-expr> ::= <unary-op> <expression>

<call-expr> ::= <expression> "(" <argument-list> ")"

<member-expr> ::= <expression> "." <identifier>
                | <expression> "[" <expression> "]"
                | <expression> "::" <identifier>

<lambda-expr> ::= "(" <parameter-list> ")" "=>" <expression>
                | "(" <parameter-list> ")" "=>" "{" <statement>* "}"

<type> ::= "number" | "string" | "bool" | "any"
         | <type> "[]"
         | "(" <type-list> ")"
         | "fn" "(" <type-list> ")" "->" <type>
\end{verbatim}
\end{definition}

\subsection{Semântica Operacional (Informal)}

\begin{definition}[Regras de Avaliação]
\textbf{Expressões Aritméticas}:
\begin{itemize}
    \item $\llbracket n_1 + n_2 \rrbracket = n_1 + n_2$ (se ambos números)
    \item $\llbracket s_1 + s_2 \rrbracket = concat(s_1, s_2)$ (se algum é string)
    \item $\llbracket n_1 / n_2 \rrbracket = n_1 / n_2$ se $n_2 \neq 0$, senão erro
\end{itemize}

\textbf{Expressões Lógicas}:
\begin{itemize}
    \item $\llbracket a \&\& b \rrbracket = \begin{cases}
        a & \text{if } isFalsy(a) \\
        b & \text{otherwise}
    \end{cases}$
    
    \item $\llbracket a || b \rrbracket = \begin{cases}
        a & \text{if } isTruthy(a) \\
        b & \text{otherwise}
    \end{cases}$
\end{itemize}

\textbf{Resolução de Variáveis}:
\begin{itemize}
    \item $\llbracket x \rrbracket_{\rho} = \rho(x)$ se $x \in dom(\rho)$
    \item $\llbracket x \rrbracket_{\rho} = \llbracket x \rrbracket_{\rho'}$ se $x \notin dom(\rho)$ e $\rho'$ é escopo pai
    \item Erro se $x$ não encontrado em nenhum escopo
\end{itemize}
\end{definition}

\section{Limitações e Considerações}

\subsection{Limitações Conhecidas}

\begin{itemize}
    \item \textbf{Performance}: Interpretador é mais lento que linguagens compiladas
    \item \textbf{Type Safety}: Sistema de tipos dinâmico não previne erros de tipo em compile-time
    \item \textbf{FFI}: Limitações ao passar tipos complexos para C
    \item \textbf{Concorrência}: Interpreter base é single-threaded (threads são isoladas)
    \item \textbf{Módulos}: Sistema de módulos sem versionamento
    \item \textbf{AOT}: Suporte experimental, nem todas features suportadas
\end{itemize}

\subsection{Características de Performance}

\begin{property}[Trade-offs]
\begin{itemize}
    \item \textbf{Adequado para}: Scripting, automação, prototipagem, aplicações pequenas/médias
    \item \textbf{Inadequado para}: Computação intensiva, sistemas real-time, aplicações críticas de performance
    \item \textbf{Bottlenecks}: Call stacks profundas, operações em arrays grandes, GC frequente
\end{itemize}
\end{property}

\section{Conclusão}

\subsection{Resumo}

Dryad é uma linguagem de programação multi-paradigma moderna que combina:
\begin{itemize}
    \item Sintaxe familiar (JavaScript/TypeScript-like)
    \item Tipagem dinâmica com anotações opcionais
    \item Modelo de execução híbrido (interpreter/bytecode/AOT)
    \item Sistema de módulos robusto com FFI
    \item Concorrência nativa (async/await e threads)
    \item Sistema de erros estruturado e informativo
\end{itemize}

\subsection{Filosofia de Design}

A linguagem prioriza:
\begin{enumerate}
    \item \textbf{Clareza e Legibilidade}: Código deve ser auto-explicativo
    \item \textbf{Flexibilidade}: Suporte a múltiplos paradigmas
    \item \textbf{Developer Experience}: Erros informativos, sintaxe consistente
    \item \textbf{Gradualidade}: Pode começar simples e adicionar complexidade quando necessário
\end{enumerate}

\subsection{Direções Futuras}

Possíveis evoluções da linguagem:
\begin{itemize}
    \item Maturação do compilador AOT
    \item Otimizações de performance no bytecode VM
    \item Expansão de módulos nativos padrão
    \item Sistema de tipos gradual mais robusto
    \item Ferramentas de análise estática
    \item Package manager integrado
\end{itemize}

\section*{Apêndice A: Tabela de Códigos de Erro}

\begin{center}
\begin{tabular}{|l|l|}
\hline
\textbf{Faixa} & \textbf{Categoria} \\
\hline
1000-1999 & Lexer Errors \\
2000-2999 & Parser Errors \\
3000-3999 & Runtime Errors \\
4000-4999 & Type Errors \\
5000-5999 & I/O Errors \\
6000-6999 & Module Errors \\
7000-7999 & Syntax Errors \\
8000-8999 & Warnings \\
9000+ & System Errors \\
\hline
\end{tabular}
\end{center}

\section*{Apêndice B: Palavras Reservadas (Completo)}

\begin{multicols}{3}
\begin{itemize}[noitemsep]
    \item async
    \item await
    \item break
    \item catch
    \item class
    \item const
    \item continue
    \item do
    \item else
    \item export
    \item extends
    \item false
    \item finally
    \item fn
    \item for
    \item from
    \item function
    \item if
    \item import
    \item in
    \item let
    \item match
    \item mutex
    \item new
    \item null
    \item private
    \item protected
    \item public
    \item return
    \item static
    \item super
    \item this
    \item thread
    \item throw
    \item true
    \item try
    \item use
    \item while
\end{itemize}
\end{multicols}

\section*{Referências}

\begin{enumerate}
    \item STANDARDIZATION\_MANIFEST.md - Dryad Project Standardization
    \item DRYAD\_LANGUAGE\_SPECIFICATION\_FOR\_LLM.md - Complete Language Specification
    \item Dryad Error Catalog - Centralized Error Definitions
    \item Dryad Source Code - Implementation Reference (Rust)
\end{enumerate}

\end{document}
